\chapter{Pleurisy (Pleuritis)}

\section{Definition}
Pleurisy is inflammation of the pleura, the double-layered membrane that surrounds the lungs and lines the chest cavity. The hallmark symptom is sharp, stabbing chest pain that worsens with breathing, coughing, or sneezing.

\section{Pathophysiology}
The pleural layers normally glide smoothly during respiration. When inflamed, they become swollen and rough, causing friction pain with each breath. Pleurisy can be dry (fibrinous) with no significant fluid accumulation, or exudative with pleural effusion (fluid between the layers). The underlying cause determines the nature of the effusion.

\section{Etiology}
\begin{itemize}
\item Viral infections (most common cause, especially coxsackievirus)
\item Bacterial pneumonia (parapneumonic effusion)
\item Pulmonary embolism with infarction
\item Tuberculosis
\item Connective tissue diseases (lupus, rheumatoid arthritis)
\item Malignancy (lung cancer, mesothelioma, metastatic disease)
\item Chest trauma or rib fracture
\item Post-cardiac injury syndrome (Dressler syndrome)
\item Asbestos exposure
\item Uremia (end-stage renal disease)
\item Pancreatitis (can cause left-sided pleural effusion)
\end{itemize}

\section{Indications for Evaluation}
\begin{itemize}
\item Sharp, stabbing chest pain that worsens with deep breathing, coughing, or sneezing
\item Pain that may radiate to the shoulder or abdomen
\item Shortness of breath (due to splinting or effusion)
\item Cough
\item Fever and chills (if infectious cause)
\item Rapid, shallow breathing
\item Friction rub heard on auscultation
\end{itemize}

\section{Related Symptoms}
\begin{itemize}
\item Chest pain (pleuritic)
\item Shortness of breath
\item Cough
\item Fever
\item Pleural effusion
\item Pneumonia
\item Pulmonary embolism
\end{itemize}

\section{Self-Care/Home Management}
\begin{itemize}
\item Take NSAIDs (ibuprofen, naproxen) for pain relief
\item Rest and avoid heavy activity
\item Apply a heating pad or ice pack to the chest
\item Use pillows to splint the chest when coughing
\item Practice deep breathing exercises to prevent atelectasis
\item Stay hydrated
\end{itemize}

\section{Diagnostic Approach}
\begin{itemize}
\item Chest X-ray to evaluate for pleural effusion or pneumonia
\item Ultrasound of the chest for pleural fluid detection
\item CT chest with contrast for detailed evaluation
\item Pleural fluid analysis (thoracentesis) for exudative effusion
\item Blood tests: CBC, ESR, CRP, D-dimer
\item ECG to rule out pericarditis or myocardial infarction
\end{itemize}

\section{Treatment/Management Overview}
\begin{itemize}
\item Treat the underlying cause (antibiotics for pneumonia, anticoagulation for PE)
\item NSAIDs or corticosteroids for pain and inflammation
\item Thoracentesis for symptomatic pleural effusion
\item Chest tube with drainage for empyema
\item Pleurodesis for recurrent malignant effusion
\item Observation for mild viral pleurisy
\end{itemize}
